\documentclass[10pt]{article}
\usepackage[T1]{fontenc}
\usepackage{lmodern,amsmath,amssymb,amsthm,microtype}
\usepackage[margin=0.8in]{geometry}
\usepackage[colorlinks=true,urlcolor=blue]{hyperref}
\newtheorem{theorem}{Theorem}
\newtheorem{lemma}{Lemma}
\newtheorem{corollary}{Corollary}
\title{Palindromic Proper-Divisor Residues at Every Modulus}
\author{Qichao Wang\\{\small Hebei University of Technology}\\[0.3em]
Daoyu Dong\\{\small University of Electronic Science and Technology of China}}
\date{October 26, 2026}
\linespread{0.96}
\hypersetup{pdftitle={Existence and Exact Finite Minima for OEIS A393833},pdfauthor={Qichao Wang and Daoyu Dong}}
\begin{document}
\maketitle
\begin{abstract}
For every $n\ge1$, we construct an integer with exactly $2n+2$
positive divisors whose $2n+1$ proper-divisor residues modulo $n$
form a palindrome. This proves the existence conjecture in OEIS
A393833. The construction needs no prime-distribution theorem:
a prime divisor of a suitable cyclotomic value suffices. An explicit
witness bound makes every minimum a finite arithmetic minimization
over ordered multiplicative partitions of $2n+2$ and prime tuples.
An exact factorization search implements that formula.
\end{abstract}

\section{An elementary witness for every index}
For $K>1$ with ordered divisors
$1=d_1<d_2<\cdots<d_D=K$, let its proper-divisor residue list be
$(d_1\bmod n,\ldots,d_{D-1}\bmod n)$. OEIS A393833 asks for the
least such $K$ with $D=2n+2$ and a palindromic list, and conjectures
that it exists at every positive index.

\begin{lemma}\label{lem:prime}
For every $n\ge2$ there is an odd prime $p\equiv1\pmod n$.
More explicitly, every prime divisor of $\Phi_n(2n)$ has these
properties, where $\Phi_n$ denotes the $n$th cyclotomic polynomial.
\end{lemma}
\begin{proof}
For $n>1$, the monic integer polynomial $\Phi_n$ has constant term
one. Its roots have absolute value one, so
\[
 \Phi_n(2n)=\prod_{\zeta\text{ primitive }n\text{th root}}(2n-\zeta)>1;
\]
indeed its absolute value is at least $(2n-1)^{\varphi(n)}>1$,
and it is positive. Choose a prime $p$ dividing this value.
No prime divisor of $2n$ divides the value, by its constant term.
Thus $p$ is odd and $p\nmid n$.
Over $\mathbf F_p$, the polynomial $X^n-1$ is square-free, since
its derivative is $nX^{n-1}$. The factorization
$X^n-1=\prod_{d\mid n}\Phi_d(X)$ therefore has pairwise coprime
factors modulo $p$. The residue of $2n$ is a root of $\Phi_n$,
so its multiplicative order is exactly $n$: an order $d<n$ would
also make it a root of a factor of $X^d-1$, contradicting that
coprimality. The order divides $p-1$, proving $p\equiv1\pmod n$.
\end{proof}

\begin{theorem}\label{thm:witness}
Let $n\ge1$, and let $p$ be any odd prime congruent to one modulo
$n$. Then
\[
 K=2p^n
\]
has exactly $2n+2$ divisors and its proper-divisor residues form a
palindrome. Thus A393833 has a positive value at every index.
\end{theorem}
\begin{proof}
Since $p>2$, its ordered divisors are
\[
 1,2,p,2p,p^2,2p^2,\ldots,p^n,2p^n.
\]
There are $2(n+1)$ of them. After deleting the final divisor,
their residues alternate $1,2,1,2,\ldots,1$ modulo $n$, a list
of odd length and hence a palindrome. For $n=1$ use $p=3$;
for $n\ge2$, Lemma~\ref{lem:prime} supplies a suitable prime.
\end{proof}

Define $p_n$ to be the least odd prime congruent to one modulo $n$
and $B_n=2p_n^n$. This is an explicit finite arithmetic construction:
$p_1=3$, and for $n\ge2$ one need test primes only up to
$\Phi_n(2n)$ by Lemma~\ref{lem:prime}. Consequently
\begin{equation}\label{eq:bound}
 a(n)\le B_n\le2\Phi_n(2n)^n\qquad(n\ge2).
\end{equation}

\section{An exact finite minimum formula}
Put $D=2n+2$. An ordered multiplicative partition of $D$ is a tuple
$(f_1,\ldots,f_t)$ with $f_i\ge2$ and $\prod_i f_i=D$;
order is retained. There are finitely many, with
$1\le t\le\lfloor\log_2 D\rfloor$. For such a tuple and primes
$q_1<\cdots<q_t$, form
\[
 M(\mathbf q,\mathbf f)=\prod_{i=1}^t q_i^{f_i-1},\qquad
 \mathcal D(\mathbf q,\mathbf f)=
 \operatorname{sort}\left\{\prod_{i=1}^t q_i^{j_i}:0\le j_i<f_i\right\}.
\]
The latter list has exactly $D$ elements; write them as
$c_1<\cdots<c_D$. Define the finite candidate set
\begin{align*}
 \mathcal C_n=\bigl\{M(\mathbf q,\mathbf f):\;&
 \prod_i f_i=D,\quad q_1<\cdots<q_t\le B_n,\\
 &M(\mathbf q,\mathbf f)\le B_n,\quad
 c_i\equiv c_{D-i}\pmod n\ (1\le i<D)\bigr\}.
\end{align*}
\begin{theorem}\label{thm:min}
For every $n\ge1$, the exact least value is
\[
 \boxed{a(n)=\min\mathcal C_n.}
\]
In particular the zero escape value never occurs. Any independently
verified witness can replace $B_n$ throughout the formula.
\end{theorem}
\begin{proof}
Theorem~\ref{thm:witness} places $B_n$ in $\mathcal C_n$, using
the ordered bases $(2,p_n)$ and factors $(2,n+1)$.
Every member has exactly $D$ divisors, and the displayed congruences
are exactly the palindrome condition on the proper-divisor list.
Conversely a least admissible integer is at most $B_n$.
Its unique prime factorization has distinct increasing bases $q_i$
and positive exponents $f_i-1$ with $\prod_i f_i=D$; all its bases
are at most $B_n$. Thus it belongs to the candidate set.
\end{proof}

This formula is a finite arithmetic minimization, not a claim of an
elementary closed expression without a minimum. The substantive new
global input is the explicit witness and bound, which removes the
unbounded existence search in the original definition.

\section{Certified pruning and implementation}
Two elementary facts greatly shorten the factor-pattern search.
First, any admissible $K$ satisfies
\begin{equation}\label{eq:necessary}
 K/q_1\equiv1\pmod n,
\end{equation}
because its largest proper divisor must have the same residue as
its first proper divisor, one. Hence any partial factorization $V$
whose smallest prime is $q_1$ can be discarded as soon as
$\gcd(V/q_1,n)>1$; multiplication cannot repair this obstruction.

Second, define $E(R)=\sum_{\ell^a\parallel R}a(\ell-1)$ for
prime $\ell$. For any multiplicative partition $R=\prod f_i$,
\begin{equation}\label{eq:exponents}
 \sum_i(f_i-1)\ge E(R).
\end{equation}
Indeed $(uv-1)\ge(u-1)+(v-1)$ for integers $u,v\ge2$;
splitting every factor into primes proves the inequality.
At a partial factorization $V$, choose the next exponent $e=f-1$,
where $f\mid R$, $f>1$, and the next prime $q$. Every completion
then has value at least $Vq^{e+E(R/f)}$. With current verified
upper bound $H$, one need enumerate only
\[
 q\le\left\lfloor\bigl((H-1)/V\bigr)^{1/(e+E(R/f))}\right\rfloor.
\]
All roots and powers in the evaluator use integer arithmetic.
It enumerates every divisor $f>1$ of the remaining divisor count,
so it retains all ordered exponent patterns, including those whose
exponents increase with their bases. It applies \eqref{eq:necessary}
before the final sorted-divisor palindrome check. Validated published
terms may initialize $H$; they serve only as witnesses, and every
smaller candidate is still excluded by the complete search.

The independent verification uses a divisor sieve through $K=200000$,
compares all minima it encounters with the factor-pattern evaluator,
checks all 13 terms stated in the official entry and its comments,
and checks the explicit witnesses through $n=100$. A b-file gives
exact computed minima through $n=12$, including an index absent
from the entry's displayed data. These finite checks supplement the
proofs; the proofs establish existence and the finite formula at
every positive index.

\begin{thebibliography}{9}
\bibitem{oeis} The OEIS Foundation Inc.,
\href{https://oeis.org/A393833}{A393833}, accessed October 26, 2026.
\end{thebibliography}
\end{document}
